%%%%%%%%%%%%%%%%%%%%%%%%%%%%%%%%%%%%%%%%%%%%%%%%%%%%%%%%%%%%%%%%%%%%%%
%%%%%%%%% Select one of the options, and comment the rest of them

%%%%%%%%%% Option 1:  to compile with pdflatex : parameter "t" - to align to the top
\documentclass[professionalfonts,t]{beamer}
%sans font?

%%%%%%%%%% Option 3: to create handout for print
%\documentclass[t,handout]{beamer}
%\usepackage{pgfpages}              % to put several slides on one page
%\pgfpagesuselayout{2 on 1}[a4paper, border shrink=5mm]             % 2 slides on 1 page
%\pgfpagesuselayout{4 on 1}[a4paper,landscape, border shrink=5mm]   % 4 slides on 1 page, and landscaped


%%%%%%%%%%%%%%%%%%%%%%%%%%%%%%%%%%%%%%%%%%%%%%%%%%%%%%%%%%%%%%%%%%%%
%%%%%%%%%%%%%% Select the Theme %%%%%%%%%%%%%%%%%%%%%%%%%%%%%%%%%%%
\usetheme{Dresden}     % OK
%\usetheme{Berlin}
%\usetheme{Bergen}      % NO
%\usetheme{Boadilla}    % NO
%\usetheme{Copenhagen}  % NO
%\usetheme{Hannover}    % NO
%\usetheme{Luebeck}     % NO
%\usetheme{Marburg}     % NO
%\usetheme{Pittsburgh}  % NO
%\usetheme{default}
%\usetheme{Singapore}   % OK
%\usetheme{boxes}
%\usecolortheme{structure}
%\usecolortheme{rose}
%\usecolortheme{beaver}


\definecolor{mymaroon}{cmyk}{0.0, 1.0, 1.0, 0.498}
\definecolor{myblue}{cmyk}{1.0, 1, 0, 0.5}
\definecolor{mygreen}{cmyk}{100, 0, 100, 50}
\setbeamercolor*{palette secondary}{use=structure,fg=white,bg=myblue}
\setbeamercolor*{palette tertiary}{use=structure,fg=white,bg=mymaroon}

%\usepackage{beamerthemesplit}              %
\beamertemplateballitem % fancy bullets and numbering

\setbeamertemplate{navigation symbols}{}   % suppress navigation symbols
\addtobeamertemplate{frametitle}{}{%
	\logo{../images/IIT_logo}
	\iffalse
	
	\begin{tikzpicture}[remember picture,overlay]
	\node[anchor=center, yshift=-13pt, xshift=-5pt] at (current page.north) 
	{\includegraphics[height=1.1cm]{../images/Argonne_cmyk_black-eps-converted-to}\hspace{10cm}};
	
	\node[anchor=north east, yshift=3pt, xshift=0pt] at (current page.north east) 
	{\includegraphics[height=0.7cm]{../images/IIT_Logo_blk}};
	\end{tikzpicture}
    
     \fi
}
% other possibilities to include LOGO. it puts it in RLC

%
%\pgfdeclareimage[width=1cm]{logo}{../images/IIT_Logo}
%\logo{\pgfuseimage{logo}}


% load additional packages

\usepackage{xcolor}
\usepackage{graphicx}
\usepackage{amsmath}
\usepackage{amssymb}
\usepackage{amsthm}
\usepackage{graphicx}
\usepackage{url}
\usepackage{color}
\usepackage{booktabs} % Allows the use of \toprule, \midrule and \bottomrule in tables
\usepackage{pifont}% http://ctan.org/pkg/pifont
\usepackage{epstopdf}
\usepackage[export]{adjustbox}
\usepackage{tikz}
\usetikzlibrary{shapes.misc}
\usetikzlibrary{shapes,arrows,decorations.markings,shadows,positioning}

% Your Abbreviations
\newcommand\bE{{\mathbb{E}}}
\newcommand\bR{{\mathbb{R}}}
\newcommand\bH{{\mathbf{H}}}
% End abbreviations

\newcommand\Wider[2][3em]{%
	\makebox[\linewidth][c]{%
		\begin{minipage}{\dimexpr\textwidth+#1\relax}
			\raggedright#2
		\end{minipage}%
	}%\textbf{}
}

%%%%%%%%%%%%%%%%%%%%% to edit the main text below
%NOTES ON SOME TECHNICS
%%%% Box %%%%%%%%%%%%%%%%%%%%%%%%%%%%%%%%%%%%%%%%%%%%%%%
%{\fbox{ \parbox[t]{10cm}{ SOME TEXT }}}

%%% include a picture. The file should be with extention EPS, e.g. FILENAME.EPS
%\begin{figure}[h]
%\centering
%\includegraphics[width=.7\linewidth]{FILENAME}
%\caption{{\footnotesize PUT_CAPTION }}
%\end{figure}

%\subtitle{}
%\institute[ANL/IIT]{Argonne National Laboratory\\Illinois Institute of Technology}

\title[]{Optimization and \\ Surrogate Models Work}
\author[N.Neveu]{{\Large Nicole Neveu}}
\institute[ANL, IIT] % (optional, but mostly needed)
{   Illinois Institute of Technology \\
	Argonne National Laboratory \\
    %\url{nneveu@anl.gov} 
}
% - Use the \inst command only if there are several affiliations.
% - Keep it simple, no one is interested in your street address.
\date{ \today \\
\includegraphics[width=3cm,keepaspectratio]{../logos/Argonne_cmyk_black}%
\hfill \hfill \hfill%
\includegraphics[width=4cm,keepaspectratio]{../logos/IIT_Logo_blk-eps-converted-to}%
}

%\date[IIT, April 2009]{
%           Space Charge 2017 \\ Oc 18, 2009  }



\begin{document}
\section{Intro}

\begin{frame}
  \titlepage
\end{frame}
\begin{frame}
	\frametitle{Outline}
	\tableofcontents
\end{frame}

% Structuring a talk is a difficult task and the following structure
% may not be suitable. Here are some rules that apply for this
% solution:

% - Exactly two or three sections (other than the summary).
% - At *most* three subsections per section.
% - Talk about 30s to 2min per frame. So there should be between about
%   15 and 30 frames, all told.

% - A conference audience is likely to know very little of what you
%   are going to talk about. So *simplify*!
% - In a 20min talk, getting the main ideas across is hard
%   enough. Leave out details, even if it means being less precise than
%   you think necessary.
% - If you omit details that are vital to the proof/implementation,
%   just say so once. Everybody will be happy with that.

\section{TBA Optimization Work}
%%%%%%%%%%%%%%%%%%%%%%%%%%%%%%%%%%%%%%%%%%%%%%%%%%%%%%%%%%%%%%%%%%%%%%%%%%%%%%%%
\begin{frame}
	\frametitle{TBA Optimization}
	Goal 1: Finish this arxiv paper in progress: \url{https://arxiv.org/pdf/1302.2889.pdf}
		\begin{itemize}
			\item Use TBA beam line as model (up to the septum)
			\item Results on next two slides 
			\item I will add AWA section to paper
			\item Hopefully submit in a few weeks, given few rounds of editing
		\end{itemize} 
	\vspace{1em}
	Goal 2: Finish whole TBA optimization
		\begin{itemize}
			\item Framework in place and working to do this
			\item Needs a few more big runs on Bebop 
			\item Will be done by/before IPAC
		\end{itemize}

\end{frame}

\begin{frame}[t]
\frametitle{Pareto front at entrance of kicker, z = 16.5 m}
		\begin{minipage}{0.48\textwidth}
			Ex-1
			\includegraphics[width=1.0\linewidth]{/home/nicole/Documents/opal_stuff/awa-optim/ex-1_gen96}
			Ex-2	\includegraphics[width=1.0\linewidth]{/home/nicole/Documents/opal_stuff/awa-optim/ex-2_gen81}
		\end{minipage}
		\begin{minipage}{0.48\textwidth}
			Ex-3
			\includegraphics[width=1.0\linewidth]{/home/nicole/Documents/opal_stuff/awa-optim/ex-3_gen53}
			Ex-4
			\includegraphics[width=1.0\linewidth]{/home/nicole/Documents/opal_stuff/awa-optim/ex-4_gen95}
		\end{minipage}

\end{frame}

\begin{frame}[t]
	\frametitle{Pareto front at entrance of septum, z = 18.5 m}
	\begin{minipage}{0.48\textwidth}
		Ex-1
		\includegraphics[width=1.0\linewidth]{/home/nicole/Documents/opal_stuff/awa-optim/ex-1_gen96_z2}
		Ex-2
		\includegraphics[width=1.0\linewidth]{/home/nicole/Documents/opal_stuff/awa-optim/ex-2_gen81_z2}
	\end{minipage}
	\begin{minipage}{0.48\textwidth}
		Ex-3
		\includegraphics[width=1.0\linewidth]{/home/nicole/Documents/opal_stuff/awa-optim/ex-3_gen53_z2}
		Ex-4
		\includegraphics[width=1.0\linewidth]{/home/nicole/Documents/opal_stuff/awa-optim/ex-4_gen95_z2}
	\end{minipage}	
\end{frame}
%%%%%%%%%%%%%%%%%%%%%%%%%%%%%%%%%%%%%%%%%%%%%%%%%%%%%%%%%%%%%%%%%%%%%%%%%%%%%%%%

%%%%%%%%%%%%%%%%%%%%%%%%%%%%%%%%%%%%%%%%%%%%%%%%%%%%%%%%%%%%%%%%%%%%%%%%%%%%%%%%
\section{Surrogate Models}
\begin{frame}
	\frametitle{Surrogate Models and Machine Learning}
	Andreas and Auralee invited me to collaboration on a Surrogate Model / Machine learning paper that is an extension of the work we did for the ML workshop at SLAC.
	\vspace{1em}	
	
	Goal 1: Generate 3 sets of data to be used for a Surrogate Model / Machine Learning.
\end{frame}


\begin{frame}[t]
	\frametitle{Pareto front, z = 18.5 m}
	\begin{minipage}{0.48\textwidth}
		Ex-1
		\includegraphics[width=1.0\linewidth]{/home/nicole/Documents/opal_stuff/surrogate-ga/0100}
		Ex-2
		\includegraphics[width=1.0\linewidth]{/home/nicole/Documents/opal_stuff/awa-optim/ex-2_gen81_z2}
	\end{minipage}
\end{frame}
\end{document}
















